\begin{table*}[t!]
  \centering
  \caption{\textbf{Experimental carbon-source growth and diagnostic enzyme phenotype prediction in BacDive.} Inputs are the binary gapseq, ModelSEED, and CarveMe calls of the gapseq comparison tables \citep{zimmermann2021gapseq}. Track 1 pools these calls with a genus-peer consensus and an unsupervised cross-assay coupling term without phenotype labels. Supervised rows are five-fold genus-blocked cross-validations (pooled out-of-fold metrics across unseen genera) of an L1-SVM, a random forest, and a gradient-boosted decision tree (GBDT) re-implemented on the reconstruction calls after the designs of Traitar, BacDive-AI, and DiASPora, and of a GBDT on the GSPA posteriors. Track 2 reveals $30\%$ of each strain\'s assays as partial phenotypic evidence and evaluates the held-out $70\%$. The carbon-source table comprises $791$ strains and $50$ substrates and the enzyme table $3{,}164$ strains after retaining assays with a definitive activity and calls from all three reconstructions. Middle rows report the 10-enzyme diagnostic panel ($50{,}075$ strain--EC pairs).}
  \label{tab:bacdive}
  \resizebox{\textwidth}{!}{%
  \begin{tabular}{llccccccc}
    \toprule
    \textbf{Assay Cohort \& Protocol} & \textbf{Baseline or Method Pipeline} & \textbf{Strains / Assays} & \textbf{Precision} & \textbf{Recall} & \textbf{$F_1$ Score} & \textbf{Bal.\ Acc.} & \textbf{ROC-AUC} & \textbf{McNemar ($+ / -$, $p$) or Runtime} \\
    \midrule
